% !TEX root = chapter-standalone.tex
\clearpage
\section{Composing relations}

\linkvideo{spring2021-relations:relations:comp-rel} % Composing relations

The visualization in \cref{fig:example_rel} hints at the fact that we can compose relations if the target of the one is the source of the other.

To illustrate composition of relations, consider a simple example, involving sets~\setA,~\setB, and~\setC, and relations~$\relA \colon \setA \sto \setB$ and~$\relB \colon \setB \sto \setC$, as depicted graphically below in \cref{fig:example_rel_composable}.

%
\begin{figure}[h!]
    \centering
    \subfloat[Relations compatible for composition. \label{fig:example_rel_composable}]{
        \includesag{30_rel_2}}\\
    \subfloat[Composition of relations. \label{fig:example_rel_composed}]{
        \includesag{30_rel_3}
    }
    \caption{Illustrations for relations composition.}
    \label{fig:rel_composition_example}
\end{figure}
%
The composite relation~$\relA \mthen \relB \colon \setA \mto \setC$ is defined to be such that~$\inrel \ela {(\relA \mthen \relB)} \elc$ if and only if there exists some~$\elb \setin \setB$ such that~$\inrel \ela \relA \elb $ and~$\inrel \elb \relB{\elc}$.
Graphically this means that for~$\tup{\ela,\elc}$ to be an element of the relation~$\relA \mthen \relB$, the elements~$\ela$ and~$\elc$ need to be connected by at least one sequence of two arrows such that the target of the first arrow is the source of the second.

For example, in \cref{fig:example_rel_composable}, there is an arrow from~$\sfondue$ to~$\swater$, and from there on to~$\sapple$, and therefore, in the composition~$\relA \mthen \relB$ depicted in~\cref{fig:example_rel_composed}, there is an arrow from~$\sfondue$ to~$\sapple$.

\begin{ctdefinition}[Relation composition]
    \label{def:composition-of-relations}
    \SYNDEF{composition of relations}
    Given relations~$\relA \colon \setA \mto \setB$,~$\relB \colon \setB \mto \setC$, their composition is the relation
    \begin{equation}
        \relA \mthen \relB \definedas \makeset{\tup{\ela,\elc} \setin \setA \cartprod \setC \mid \exists \elb \setin \setB \colon \pars{\inrel{\ela}{\relA}{\elb} } \booland \pars{\inrel{\elb}{\relB}{\elc}}}.
    \end{equation}
\end{ctdefinition}

\todotext{JL: @TAs - write up a simple, nice example of relation composition here and put it in an ``example'' environment}

\todotext{JL: @TAs - write up two simple examples of relation composition and put them here as (non-graded) exercises, plus write in the solutions using the solutions environment}

\begin{gradedexercise}[ComposingRelations]
    \
    \begin{enumerate}
        \item Let~$\setA = \natnumbers$,~$\setB = \wnumbers$, and~$\setC = \reals$.
              Consider the relation~$\mora \colon \setA \mto \setB$ with
              \begin{equation}
                  \mora = \makeset{ \tup{\ela, \elb} \setin \natnumbers \cartprod \wnumbers \mid \ela = \elb^2 },
              \end{equation}
              and consider the relation~$\morb \colon \setB \mto \setC$ with
              \begin{equation}
                  \morb = \makeset{ \tup{\elb, \elc} \setin \wnumbers \cartprod \reals \mid \elb = 2 \elc}.
              \end{equation}
              Calculate the relation~$\mora \mthen \morb \colon \setA \mto \setC$.

        \item Let~$\setA = \natnumbers$,~$\setB = \wnumbers \cartprod \wnumbers$, and~$\setC = \reals$.
              Consider the relation~$\mora \colon \setA \mto \setB$ with
              \begin{equation}
                  \mora = \makeset{\tup{\ela, \tup{\elnb{1}, \elnb{2}}} \setin \natnumbers \cartprod (\wnumbers \cartprod \wnumbers) \mid \ela = \elnb{1} - \elnb{2}}
              \end{equation}
              and consider the relation~$\morb \colon \setB \mto \setC$ with
              \begin{equation}
                  \morb = \makeset{ \tup{\tup{\elnb{1}, \elnb{2}}, \elc} \setin (\wnumbers \cartprod \wnumbers) \cartprod \reals \mid \elnb{2} \elc = \elnb{1}}.
              \end{equation}

              Calculate the relation~$\mora \mthen \morb \colon \setA \mto \setC$.
    \end{enumerate}
\end{gradedexercise}

\begin{gradedexercise}[\exname{ComposingFiniteRelations}]
    \label{ex:ComposingFiniteRelations}
    \
    \begin{enumerate}
        \item Consider finite sets~$\setA = \makeset{a_1, \dots, a_n}$,~$\setB = \makeset{b_1, \dots, b_m}$, and~$\setC = \makeset{c_1, \dots, c_p}$, along with relations~$\relA \colon \setA \mto \setB$ and~$\relB \colon \setB \mto \setC$.
              With each relation, we associate a matrix:
              \[
                  M_{\relA} \in \reals^{n \times m} \quad \text{and} \quad M_{\relB} \in \reals^{m \times p},
              \]
              defined as follows:
              \[
                  (M_{\relA})_{ij} = \begin{cases}
                      1, & (a_i, b_j) \setin \relA, \\
                      0, & \text{otherwise}
                  \end{cases}
              \]
              and
              \[
                  (M_{\relB})_{jk} = \begin{cases}
                      1, & (b_j, c_k) \setin \relB, \\
                      0, & \text{otherwise}
                  \end{cases}
              \]
              where \( i \in \makeset{1, \dots, n} \), \( j \in \makeset{1, \dots, m} \), and \( k \in \makeset{1, \dots, p} \).
              Similarly, for the composition we define the matrix~$M_{\relA \mthen \relB} \in \reals^{n \times p}$ by
              \[
                  (M_{\relA \mthen \relB})_{ik} = \begin{cases}
                      1, & (a_i, c_k) \setin \relA \mthen \relB, \\
                      0, & \text{otherwise}
                  \end{cases}.
              \]
              Show that
              \[
                  (M_{\relA \mthen \relB})_{ik} = \max_{j} \{ (M_{\relA})_{ij} \cdot (M_{\relB})_{jk} \}.
              \]

        \item Consider~$\setA = \setB = \setC = \makeset{1, 2, 3}$, with~$\relA = \makeset{(1,1), (2,2), (3,3)}$ and~$\relB = \makeset{(1,3), (2,1)}$.
              Compute~$M_{\relA}$,~$M_{\relB}$, and~$M_{\relA \mthen \relB}$, and write down~$\relA \mthen \relB$ as a set.
    \end{enumerate}
\end{gradedexercise}

\solutionof{ComposingFiniteRelations}

\subsection{Identity relations}

We have already discussed how, for any set $\setA$, there is always an identity function
\begin{equation}
    \defmapperiod{
        \mapidat\setA
    }{
        \setA
    }{
        \sto
    }{
        \setA
    }{
        \elb
    }{
        \elb
    }
    % (\elb) = \elb \quad \forall  \elb \setin\setB.
\end{equation}
which ``does nothing''.
If we turn such functions into relations, we call the result \emph{identity relations}.

\begin{definition}\label{def:identity-relation}
    Let $\setA$ be any set.
    The \emph{identity relation} on $\setA$ is
    \begin{equation}
        \catid_\setA = \makeset { \tup{\ela, \elb} \setin \setA \cartprod \setA \mid \ela = \elb}.
    \end{equation}
\end{definition}

\todotext{2026-07, JL: say then that the identity relations are neutral for composition}

\subsection{Composing functions}

If we define functions as special kinds of relations, how is relation composition related to the ``usual'' way of composition of functions?
The answer is that these two apparently different ways of composing functions actually give the same result.

\begin{lemma}
    \label{lem:comprelfun}
    Let~$\relA \setsubseteq \setA \cartprod \setB$ and~$\relB \setsubseteq \setB \cartprod \setC$ be relations which are functions.
    Then their composition~$\relA \mthen \relB \setsubseteq \setA \cartprod \setC$ is again a function, and it corresponds to the ``usual'' composition of the functions corresponding to \relA and $\relB$.
\end{lemma}

\begin{proof}
    First let us check that when \relA and \relB are composed as relations, the result is again a function.
    For this we check that~$\relA \mthen \relB$ satisfies the two conditions stated in \cref{def:functions_as_relations}.

    \begin{enumerate}
        \item Choose an arbitrary~$\ela \setin \setA$.
              We need to show that there exists~$\elc \setin \setC$ such that~$\inrel{\ela}{\relA \mthen \relB}{\elc}$.
              Since~\relA is a function, there exists~$\elb \setin \setB$ such that~$\inrel{\ela}{\relA}{\elb}$.
              Choose such a~$\elb \setin \setB$.
              Then, because~\relB is a function, there exists~$\elc \setin \setC$ such that~$\inrel{\elb}{\relB}{\elc}$.
              By the definition of composition of relations, we see that~$\elc$ is such that~$\inrel{\ela}{\relA \mthen \relB}{\elc}$.
        \item Let~$\inrel{\elna{1}}{\relA \mthen \relB}{\elnc{1}}$,~$\inrel{\elna{2}}{\relA \mthen \relB}{\elnc{2}}$.
              We need to show that if~$\ela_1 = \ela_2$, then~$\elc_1 = \elc_2$.
              So suppose~$\elna{1} = \elna{2}$.
              Since~$\inrel{\elna{1}}{\relA \mthen \relB}{\elnc{1}}$,~$\inrel{\elna{2}}{\relA \mthen \relB}{\elnc{2}}$, there exist~$\elnb{1}, \elnb{2} \setin \setB$ such that, respectively,
              \begin{equation}
                  \inrel{\elna{1}}{\relA}{\elnb{1}} \booland \inrel{\elnb{1}}{\relB}{\elnc{1}},
              \end{equation}
              \begin{equation}
                  \inrel{\elna{2}}{\relA}{\elnb{2}} \booland \inrel{\elnb{2}}{\relB}{\elnc{2}}.
              \end{equation}
              Since~$\elna{1} = \elna{2}$ and~\relA is a function, we conclude that~$\elnb{1} = \elnb{2}$ must hold.
              Now, since~\relB is also a function, this implies that~$\elnc{1} = \elnc{2}$, which is what was to be shown.
    \end{enumerate}

    Second, let us check that relation composition of functions gives the same result as the ``usual'' composition of functions.
    Let~$\mapa_{\relA}$ and~$\mapb_{\relB}$ denote the relations~\relA and~\relB when we are thinking of them in the ``usual'' way of thinking about functions.
    Our goal is to show that~$\mapa_{\relA} \mthen \mapb_{\relB}$ corresponds to~$\relA \mthen \relB$; in other words, that the latter is the graph of former.

    Suppose first that~$\tup{\ela, \elc}$ is in the graph of~$\mapa_{\relA} \mthen \mapb_{\relB}$, so~$\elc = (\mapa_{\relA} \mthen \mapb_{\relB})(\ela)$.
    In particular~$\elc = \mapb_{\relB}(\mapa_{\relA}(\ela))$, which means there exists~$\elb = \mapa_{\relA}(\ela) \setin \setB$ such that~$\tup{\ela, \elb} \setin \relA$ and~$\tup{\elb, \elc} \setin \relB$.
    This implies that~$\tup{\ela, \elc} \setin \relA \mthen \relB$.

    Conversely, suppose~$\tup{\ela, \elc} \setin \relA \mthen \relB$.
    By the definition of relation composition there must exist~$\elb \setin \setB$ such that~$\tup{\ela, \elb} \setin \relA$ and~$\tup{\elb, \elc} \setin \relB$, which means~$\elb = \mapa_{\relA}(\ela)$ and~$\elc = \mapb_{\relB}(\elb)$.
    Thus,~$\elc = \mapb_{\relB}(\mapa_{\relA}(\ela))$.
\end{proof}

